

\subsection{Introduction: System Overview}
The proposed smart healthcare robot is a mobile, autonomous platform designed to assist hospital staff in routine, non-critical tasks while providing patients with an intuitive human--robot interaction (HRI) interface. The robot operates primarily in indoor clinical environments such as hospital wards, corridors, and waiting areas. Its main roles include autonomous navigation, safe obstacle avoidance, point-to-point delivery of small items, basic patient interaction via a touch interface, and remote telepresence capabilities for clinical staff.

From a functional perspective, the robot must be able to localize itself in the hospital environment, plan safe paths between waypoints, and react to dynamic obstacles such as people, trolleys, and doors. It should expose a user-friendly interface on a tablet for nurses, doctors, and patients to issue tasks or receive information. Optional features include autonomous docking at a charging station to support extended operation and reduced manual maintenance. The system is conceived as a modular platform so that sensors, software modules, and interaction modalities can be extended in future iterations without major redesign.

The overall architecture follows a layered approach:
\begin{itemize}
    \item A \textbf{mobile base and mechanical structure} providing locomotion, payload mounting, and physical robustness.
    \item A \textbf{computing and control layer} that integrates a high-level computer running ROS~2 for navigation and coordination, and a microcontroller for real-time motor and sensor interfacing.
    \item A \textbf{navigation and perception layer} Robot able to localize and navigate in indooe environments using one or multiple navigation approaches (e.g., LiDAR-based SLAM, waypoints following, or sensor fusion combinations) with selectable sensors such as LiDAR, IMU, wheel encoders, and auxiliary safety sensors for localization, mapping, and obstacle detection.
    \item A \textbf{communication and HRI layer} offering a tablet-based user interface, web dashboard, and wireless connectivity for remote monitoring and teleoperation.
    \item A \textbf{power subsystem} supplying stable power to all components, with optional docking capabilities for automatic recharging through a docking station.
\end{itemize}

This conceptual design aims to balance technical performance, patient safety, modularity, and cost, providing a realistic path from prototype to potential deployment in Omani hospitals.

\subsection{Design Alternatives and Selection}
To arrive at a suitable system architecture for the smart healthcare robot, several design alternatives were considered. These alternatives focus on the choice of computing platform, sensor suite, mobility base, and interaction modalities, under constraints of cost, development time, and hospital safety requirements. Three main alternatives were initially evaluated:

\subsubsection{Alternative 1 -- (Reserved)}
% Intentionally left blank for future description of Alternative 1.

\subsubsection{Alternative 2 -- Selected Modular Service Robot Platform}
The selected conceptual design is a modular indoor service robot platform built around a differential or 4WD mecanum/omni drive base, Then integrated a Raspberry~Pi~5 as the main computer running ROS~2, and an ESP32 microcontroller handling low-level motor and sensor control. This design emphasizes the use of widely available, well-supported components to accelerate development while ensuring flexibility for future upgrades.

\paragraph{Architecture and Main Components}
The system is organized into six main subsystems.

\textbf{1. Computing and Control}
\begin{itemize}
    \item \textbf{Raspberry Pi 5 (4~GB)}: Serves as the main onboard computer. It runs Ubuntu with ROS~2, including SLAM and the Nav2 navigation stack. The Pi is responsible for high-level decision making, global and local path planning, sensor fusion at the navigation level, and communication with the user interfaces.
    \item \textbf{ESP32 development board}: Used as a dedicated microcontroller for real-time tasks such as reading low-level sensors (wheel encoders, IMU, some safety sensors) and generating motor control signals via the motor driver. It communicates with the Raspberry Pi over a serial or Wi-Fi link, exposing a simple command interface (e.g., velocity commands and feedback messages).
    \item \textbf{Motor driver board (12~V, 4-channel)}: Interfaces the ESP32 with the four DC motors of the 4WD base, providing the required current and voltage switching. It enables forward/backward and individual wheel control, supporting holonomic motion with omni or mecanum wheels.
    \item \textbf{Logic-level converter}: Used when necessary to safely interface 3.3~V logic (ESP32, Raspberry Pi GPIO) with any 5~V-compatible sensors or driver inputs.
\end{itemize}

\textbf{2. Navigation and Sensors}
\begin{itemize}
    \item \textbf{RPLIDAR A1}: A 360$^\circ$ LiDAR sensor providing range measurements around the robot. It is the primary sensor for 2D SLAM, map building, and obstacle detection in indoor corridors and rooms.
    \item \textbf{MPU6050 IMU}: Combines a 3-axis accelerometer and 3-axis gyroscope to provide inertial measurements used for pose estimation and to improve localization robustness, especially during short LiDAR occlusions.
    \item \textbf{Wheel encoders (\(4\times\))}: Provide odometry information by measuring wheel rotations. Odometry is fused with LiDAR and IMU data to estimate the robot pose in real time.
    \item \textbf{Front camera}: A forward-facing camera supports teleoperation, visual monitoring, and future features such as vision-based docking or patient recognition.
    \item \textbf{Safety sensors}: Ultrasonic and infrared time-of-flight (ToF) sensors provide redundant short-range obstacle detection around the robot, particularly at low heights where LiDAR coverage may be limited. Optional bumper switches add a final physical safety layer in case of contact. An emergency stop button allows staff to immediately cut motor power in abnormal situations.
\end{itemize}

\textbf{3. Communication and User Interface}
\begin{itemize}
    \item \textbf{Android tablet}: Mounted on the robot as the main HRI screen. It can display a custom application for task selection, patient interaction (e.g., messages, questionnaires), and robot status.
    \item \textbf{Web dashboard}: Accessible via the hospital Wi-Fi network to allow supervisors to monitor robot status, view maps, send navigation goals, or schedule tasks from a desktop browser.
    \item \textbf{Mobile application}: A lightweight mobile app (or responsive web interface) provides remote control or status monitoring for authorized staff when they are not near the robot.
    \item \textbf{Wi-Fi connectivity}: The onboard Raspberry Pi uses its built-in Wi-Fi module to connect to the hospital network or a dedicated access point, enabling communication with the tablet, web dashboard, and cloud services if required.
\end{itemize}

\textbf{4. Robot Base and Mechanical Structure}
\begin{itemize}
    \item \textbf{4WD omni/mecanum base}: A four-wheel drive base with omni or mecanum wheels provides high maneuverability in narrow hospital corridors, allowing sideways motion and precise docking alignment.
    \item \textbf{DC motors with gearboxes}: Provide sufficient torque at low speeds for safe indoor motion and payload carrying (e.g., small items, medical supplies).
    \item \textbf{Aluminum extrusion frame}: A modular aluminum profile structure supports the tablet mount, electronics plates, LiDAR, and sensors. It simplifies mechanical integration and future modifications.
    \item \textbf{Mounting and cabling}: A dedicated electronics mounting plate holds the Raspberry Pi, ESP32, motor driver, power converters, and wiring. Tablet mounts, top plates, brackets, and cable channels (with cable ties) ensure robustness, maintainability, and neat cable management.
\end{itemize}

\textbf{5. Power System}
\begin{itemize}
    \item \textbf{12~V LiFePO$_4$ battery (10--20~Ah)}: Provides the main energy source for the robot with higher safety and cycle life compared to typical Li-ion packs. The capacity is chosen to support several hours of operation between charges, depending on duty cycle.
    \item \textbf{Battery Management System (BMS)}: Monitors cell voltages and currents, providing overcharge, over-discharge, and short-circuit protection to improve safety and battery lifespan.
    \item \textbf{DC--DC buck converters}: One 12~V to 5~V 5~A converter powers the Raspberry Pi, LiDAR, and camera, while a second 12~V to 5~V 2~A converter powers the ESP32 and low-power sensors. This separation improves noise isolation and reliability.
    \item \textbf{Protective elements}: An inline fuse (around 10~A), master power switch, XT60 connectors, and an optional power distribution board simplify safe power routing and make maintenance easier.
\end{itemize}

\textbf{6. Automatic Docking System}
\begin{itemize}
    \item \textbf{Robot-side hardware}: Magnetic or spring-loaded charging contacts mounted at the rear or underside of the robot provide reliable electrical connection when docked. Optional visual (ArUco markers) or IR beacons assist the robot in aligning precisely with the dock.
    \item \textbf{Docking station}: Consists of two metal charging plates connected to a 12~V, 5~A AC--DC adapter, integrated into a mechanical cradle that guides the robot into the correct position. An LED status indicator can show charging state for quick visual feedback.
\end{itemize}

\paragraph{Advantages}
The chosen design offers several advantages:
\begin{itemize}
    \item Use of Raspberry Pi 5 and ESP32 leverages strong community support, ROS~2 compatibility, and abundant libraries, reducing development risk.
    \item LiFePO$_4$ battery chemistry, BMS, and clear power distribution improve electrical safety, which is critical in hospital environments.
    \item Modular mechanical and electronic structure allows incremental upgrades (e.g., better sensors, additional HRI modules) without redesigning the entire robot.
    \item LiDAR-based SLAM combined with IMU and odometry provides robust indoor navigation in dynamic hospital spaces.
    \item Tablet and web dashboard interfaces improve usability for both patients and staff, supporting HRI and remote supervision.
\end{itemize}

\paragraph{Limitations}
Despite its strengths, this alternative has some limitations:
\begin{itemize}
    \item The Raspberry Pi is less powerful than industrial PCs, which may limit the complexity of perception or AI algorithms that can run onboard.
    \item LiDAR-based navigation can struggle with glass surfaces, highly reflective materials, or crowded environments, requiring careful tuning and safety margins.
    \item The 4WD mecanum/omni base is mechanically and control-wise more complex than a simple differential drive, potentially increasing integration and maintenance effort.
    \item Automatic docking requires precise mechanical alignment and robust perception; misalignment can reduce charging reliability and may require fallback manual docking procedures.
    \item As a prototype platform, the design may require further hardening (e.g., ingress protection, cable shielding) before being deployed in real hospital wards.
\end{itemize}

Overall, this alternative provides a balanced compromise between performance, cost, safety, and development complexity, making it a suitable conceptual design for the proposed smart healthcare robot.






\clearpage